\documentclass[12pt]{article}

% Idioma e codificação
\usepackage[brazil]{babel}
\usepackage[utf8]{inputenc}
\usepackage[T1]{fontenc}

% Matemática e listas
\usepackage{amsmath, amssymb, amsfonts}
\usepackage{enumitem}
\usepackage{tabularx}
\usepackage{booktabs}

% Layout e cores
\usepackage{geometry}
\usepackage{graphicx}
\usepackage{xcolor}

% Cores institucionais
\definecolor{maincolor}{RGB}{0,100,50}
\definecolor{forestgreen}{RGB}{0,100,50}

% Cabeçalho e rodapé
\usepackage{fancyhdr}

% TikZ configuration for diagrams
\usepackage{tikz}
\usetikzlibrary{shapes.geometric, shapes.multipart, arrows, positioning, calc}

% Blocos coloridos
\usepackage[most]{tcolorbox}
\newtcolorbox{block}[1]{%
  colback=green!5,
  colframe=maincolor,
  coltitle=white,
  title=#1,
  fonttitle=\bfseries,
  boxrule=1.5pt,
  arc=3mm,
  left=6pt, right=6pt, top=6pt, bottom=6pt
}

% Espaço para resposta
\newtcolorbox{resposta}{
  colback=white,
  colframe=gray!40,
  arc=0mm,
  height=8cm,
  left=6pt, right=6pt, top=6pt, bottom=6pt,
  title=\textit{Resposta:},
  coltitle=gray!80,
  fonttitle=\bfseries\small,
  attach title to upper
}

% Margens
\geometry{margin=2.5cm}

% Fancyhdr
\pagestyle{fancy}
\fancyhf{}

\renewcommand{\headrulewidth}{2pt}
\renewcommand{\footrulewidth}{2pt}

\renewcommand{\headrule}{\hbox to\headwidth{%
  \color{forestgreen}\leaders\hrule height \headrulewidth\hfill}}

\renewcommand{\footrule}{\hbox to\headwidth{%
  \color{forestgreen}\leaders\hrule height \footrulewidth\hfill}}

\setlength{\headheight}{52pt}
\addtolength{\topmargin}{-20pt}

\fancyhead[L]{%
  \includegraphics[height=1.2cm]{Universidade_de_Rio_Verde.jpg}
}
\fancyhead[R]{%
  \raggedleft
  \textcolor{forestgreen}{%
    \small
    \textbf{Lista de Exercícios}\\
    Me.~Sergio Souza Novak\\
    \textit{Banco de Dados}
  }
}
\fancyfoot[C]{\textcolor{forestgreen}{\thepage}}

\begin{document}

\vspace*{0.2cm}
\begin{center}
  {\Large \textbf{Lista de Exercícios: Normalização}}\\[0.2cm]
  {\large Banco de Dados}\\[0.2cm]
  Tema: 4ª Forma Normal (4FN) --- Dependências Multivaloradas
\end{center}

\vspace{0.3cm}

\begin{block}{Instruções}
Para cada questão: (1) identifique as Dependências Multivaloradas (DMV) e Funcionais (DF); (2) aplique o teste da 4FN verificando se o determinante é superchave; (3) apresente a decomposição se necessário, indicando PKs e FKs.
\end{block}

\vspace{0.4cm}

\begin{enumerate}[label=\textbf{Questão \arabic*:}, leftmargin=*]

% ================================================================
%  BLOCO 1 — CONCEITUAL
% ================================================================
  \item \textbf{Qual das seguintes situações tipicamente sinaliza a necessidade de aplicar a 4FN?}
  \begin{enumerate}[label=\alph*)]
    \item Presença de atributos derivados (ex: Total = Preço * Qtd).
    \item Tabelas com chaves compostas por muitos atributos.
    \item Atributos que aceitam valores Nulos.
    \item Existência de "fatos independentes muitos-para-muitos" sobre uma entidade.
    \item Um atributo que depende apenas de parte da chave composta.
  \end{enumerate}

  \item \textbf{O uso da 4FN é crucial em sistemas OLTP (Online Transaction Processing) principalmente para:}
  \begin{enumerate}[label=\alph*)]
    \item Facilitar a geração de relatórios complexos.
    \item Garantir que o banco de dados ocupe o máximo de espaço possível.
    \item Minimizar anomalias durante operações de INSERT, UPDATE e DELETE.
    \item Evitar o uso de chaves estrangeiras.
    \item Aumentar a velocidade de buscas envolvendo produtos cartesianos.
  \end{enumerate}

  \item \textbf{(Conceito)} Explique, com suas palavras, o que é uma \textbf{Dependência Multivalorada (DMV)} e qual a principal diferença entre ela e uma \textbf{Dependência Funcional (DF)}. Dê um exemplo simples de cada.
  \begin{resposta}
  \end{resposta}

  \vspace{0.4cm}

  \item \textbf{(Definição)} Quais são os dois requisitos fundamentais para que uma relação seja considerada na \textbf{4ª Forma Normal (4FN)}? Por que uma relação deve obrigatoriamente estar na BCNF antes de ser testada para a 4FN?
  \begin{resposta}
  \end{resposta}

  \vspace{0.4cm}

  \item \textbf{(Redundância)} Por que a presença de duas ou mais DMVs independentes em uma mesma tabela gera o chamado ``Produto Cartesiano''? Quais são os problemas práticos (anomalias) que esse fenômeno causa no banco de dados?
  
  \vspace{0.2cm}
  \begin{tikzpicture}[node distance=2cm, font=\small\sffamily]
      \node (x) [draw, circle, fill=orange!20, thick] {X};
      \node (y) [draw, rectangle, right=of x, fill=blue!10] {\{Y1, Y2, \dots\}};
      \node (z) [draw, rectangle, left=of x, fill=green!10] {\{Z1, Z2, \dots\}};

      \draw[->>, thick, blue] (x) -- (y) node[midway, above] {\tiny DMV};
      \draw[->>, thick, green!60!black] (x) -- (z) node[midway, above] {\tiny DMV};

      \node [below=0.5cm of x] {\small $R(x) = Y(x) \times Z(x)$};
  \end{tikzpicture}
  % 
  
  \begin{resposta}
  \end{resposta}
  
  \item \textbf{Em um cenário de "Projetos, Ferramentas e Engenheiros", se um projeto usa várias ferramentas e possui vários engenheiros independentes, a violação da 4FN exigiria:}
  \begin{enumerate}[label=\alph*)]
    \item Que cada engenheiro fosse associado a apenas uma ferramenta.
    \item Que cada ferramenta fosse usada por apenas um engenheiro.
    \item Que cada combinação de (engenheiro, ferramenta) aparecesse em uma linha distinta para o mesmo projeto.
    \item Que a chave primária fosse apenas o código do Engenheiro.
    \item Nenhuma das anteriores.
  \end{enumerate}

  \item \textbf{A decomposição para a 4FN é dita "sem perdas" (lossless join decomposition) porque:}
  \begin{enumerate}[label=\alph*)]
    \item Podemos reconstruir a relação original através de uma junção natural das tabelas decompostas.
    \item Não se perde espaço em disco.
    \item Todas as chaves estrangeiras são removidas.
    \item O número de linhas permanece o mesmo.
    \item Ela apaga dados irrelevantes.
  \end{enumerate}

  \item \textbf{Se uma relação possui apenas dois atributos $R(A, B)$, ela:}
  \begin{enumerate}[label=\alph*)]
    \item Pode violar a 4FN.
    \item Está obrigatoriamente na 4FN, pois não pode haver DMV não trivial.
    \item Viola a 1FN por ser pequena demais.
    \item Necessita de um terceiro atributo para ser normalizada.
    \item Sempre terá redundância exponencial.
  \end{enumerate}
% ================================================================
%  BLOCO 2 — TESTES E IDENTIFICAÇÃO
% ================================================================

  \vspace{0.4cm}

  \item \textbf{(Teste 4FN -- Pesquisadores)} A tabela abaixo registra os projetos e os idiomas de cada pesquisador. Sabendo que projetos e idiomas são independentes:
  
  \vspace{0.2cm}
  \small
  \begin{tabular}{|l|l|l|}
    \hline
    \textbf{Pesquisador (PK)} & \textbf{Projeto (PK)} & \textbf{Idioma (PK)} \\ \hline
    Ana Souza  & IA Climática   & Inglês    \\ \hline
    Ana Souza  & IA Climática   & Espanhol  \\ \hline
    Ana Souza  & IA Climática   & Francês   \\ \hline
    Ana Souza  & Robótica       & Inglês    \\ \hline
    Ana Souza  & Robótica       & Espanhol  \\ \hline
    Ana Souza  & Robótica       & Francês   \\ \hline
    Ana Souza  & Bioinformática & Inglês    \\ \hline
    Ana Souza  & Bioinformática & Espanhol  \\ \hline
    Ana Souza  & Bioinformática & Francês   \\ \hline
    Bruno Lima & Robótica       & Inglês    \\ \hline
    Bruno Lima & Robótica       & Alemão    \\ \hline
    Bruno Lima & Bioinformática & Inglês    \\ \hline
    Bruno Lima & Bioinformática & Alemão    \\ \hline
    Carla Dias & IA Climática   & Português \\ \hline
    Carla Dias & IA Climática   & Inglês    \\ \hline
    Carla Dias & Robótica       & Português \\ \hline
    Carla Dias & Robótica       & Inglês    \\ \hline
    Daniel Paz & Bioinformática & Italiano  \\ \hline
    Daniel Paz & Bioinformática & Inglês    \\ \hline
    Daniel Paz & IA Climática   & Italiano  \\ \hline
  \end{tabular}
  % 
  \vspace{0.2cm}
  \begin{enumerate}[label=\alph*)]
    \item Liste as DMVs existentes.
    \item A tabela viola a 4FN? Justifique verificando a superchave.
    \item Como ficariam as tabelas após a decomposição?
  \end{enumerate}
  \begin{resposta}
  \end{resposta}

  \vspace{0.4cm}

  \item \textbf{(Teste 4FN -- Representantes)} Em uma empresa, cada \textbf{representante} atende a várias \textbf{regiões} e vende várias \textbf{linhas de produtos}. As regiões e as linhas de produtos são independentes.
  
  \vspace{0.2cm}
  \small
  \begin{tabular}{|l|l|l|}
    \hline
    \textbf{Representante (PK)} & \textbf{Região (PK)} & \textbf{Linha\_Produto (PK)} \\ \hline
    Carlos Silva  & Norte   & Ferramentas  \\ \hline
    Carlos Silva  & Norte   & Máquinas     \\ \hline
    Carlos Silva  & Norte   & EPIs         \\ \hline
    Carlos Silva  & Sul     & Ferramentas  \\ \hline
    Carlos Silva  & Sul     & Máquinas     \\ \hline
    Carlos Silva  & Sul     & EPIs         \\ \hline
    Carlos Silva  & Leste   & Ferramentas  \\ \hline
    Carlos Silva  & Leste   & Máquinas     \\ \hline
    Carlos Silva  & Leste   & EPIs         \\ \hline
    João Paulo    & Centro  & Ferramentas  \\ \hline
    João Paulo    & Centro  & Máquinas     \\ \hline
    João Paulo    & Oeste   & Ferramentas  \\ \hline
    João Paulo    & Oeste   & Máquinas     \\ \hline
    Maria Alice   & Sul     & Eletrônicos  \\ \hline
    Maria Alice   & Sul     & Papelaria    \\ \hline
    Maria Alice   & Norte   & Eletrônicos  \\ \hline
    Maria Alice   & Norte   & Papelaria    \\ \hline
    Ricardo Lima  & Centro  & Móveis       \\ \hline
    Ricardo Lima  & Centro  & Decoração    \\ \hline
    Ricardo Lima  & Centro  & Iluminação   \\ \hline
  \end{tabular}
  % 
  \vspace{0.2cm}
  \begin{enumerate}[label=\alph*)]
    \item Identifique a Chave Primária da tabela original.
    \item Demonstre a ocorrência do produto cartesiano para o representante Carlos Silva.
    \item Aplique a decomposição para 4FN.
  \end{enumerate}
  \begin{resposta}
  \end{resposta}

  \vspace{0.4cm}
  \item \textbf{A anomalia de atualização na 4FN ocorre porque:}
  \begin{enumerate}[label=\alph*)]
    \item Ao alterar um valor de um atributo multivalorado, é necessário atualizar múltiplas linhas para manter a consistência do produto cartesiano.
    \item O banco de dados bloqueia a tabela para qualquer escrita.
    \item Não se pode inserir chaves nulas.
    \item A chave primária é muito curta.
    \item Os atributos não-chave dependem de apenas parte da chave.
  \end{enumerate}

  \item \textbf{Historicamente, a 4FN foi introduzida por Ronald Fagin em 1977 para lidar com falhas de qual forma normal anterior?}
  \begin{enumerate}[label=\alph*)]
    \item 1FN.
    \item 2FN.
    \item 3FN.
    \item BCNF (Forma Normal de Boyce-Codd).
    \item 5FN.
  \end{enumerate}
  \item \textbf{(Teste 4FN -- Hospital)} Considere a relação \texttt{MEDICO\_VINCULO}:
  
  \vspace{0.2cm}
  \small
  \begin{tabular}{|l|l|l|}
    \hline
    \textbf{Medico (PK)} & \textbf{Hospital (PK)} & \textbf{Convenio (PK)} \\ \hline
    Dr. Arnaldo & Sta. Casa   & Unimed  \\ \hline
    Dr. Arnaldo & Sta. Casa   & Bradesco \\ \hline
    Dr. Arnaldo & Sta. Casa   & Amil     \\ \hline
    Dr. Arnaldo & Hosp. Luz   & Unimed  \\ \hline
    Dr. Arnaldo & Hosp. Luz   & Bradesco \\ \hline
    Dr. Arnaldo & Hosp. Luz   & Amil     \\ \hline
    Dr. Arnaldo & Clin. Sul   & Unimed  \\ \hline
    Dr. Arnaldo & Clin. Sul   & Bradesco \\ \hline
    Dr. Arnaldo & Clin. Sul   & Amil     \\ \hline
    Dra. Beatriz & Sta. Casa   & Cassi    \\ \hline
    Dra. Beatriz & Sta. Casa   & Ipasgo   \\ \hline
    Dra. Beatriz & Hosp. Luz   & Cassi    \\ \hline
    Dra. Beatriz & Hosp. Luz   & Ipasgo   \\ \hline
    Dr. Claudio  & Clin. Norte & Unimed   \\ \hline
    Dr. Claudio  & Clin. Norte & Bradesco \\ \hline
    Dr. Claudio  & Clin. Norte & Cassi    \\ \hline
    Dr. Claudio  & Hosp. Luz   & Unimed   \\ \hline
    Dr. Claudio  & Hosp. Luz   & Bradesco \\ \hline
    Dr. Claudio  & Hosp. Luz   & Cassi    \\ \hline
    Dr. Claudio  & Sta. Casa   & Unimed   \\ \hline
  \end{tabular}
  
  \vspace{0.2cm}
  \begin{enumerate}[label=\alph*)]
    \item Quais são as DMVs?
    \item Se o Dr. Arnaldo passar a atender pelo convênio ``Cassi'', quantas linhas devem ser inseridas?
    \item Decompunha a tabela para eliminar a redundância.
  \end{enumerate}
  \begin{resposta}
  \end{resposta}

  \vspace{0.4cm}

  \item \textbf{(Análise BCNF vs 4FN)} Uma tabela \texttt{CURSOS} possui os atributos: \texttt{\{Curso, Livro\_Texto, Professor\}}. As regras são: (1) Para cada curso, há vários livros sugeridos; (2) Para cada curso, há vários professores que podem ministrá-lo; (3) Professores e livros são independentes.
  
  \begin{center}
  \begin{tikzpicture}[node distance=1.5cm, font=\footnotesize\sffamily]
      \node (c) [draw, fill=orange!10, rectangle, rounded corners] {Curso (Chave?)};
      \node (l) [draw, fill=blue!5, rectangle, right=of c] {Livro\_Texto};
      \node (p) [draw, fill=green!5, rectangle, left=of c] {Professor};
      
      \draw[->>, thick, blue] (c) -- (l);
      \draw[->>, thick, green!60!black] (c) -- (p);
  \end{tikzpicture}
  \end{center}

  \begin{enumerate}[label=\alph*)]
    \item Essa tabela está na BCNF? Justifique.
    \item Essa tabela está na 4FN? Justifique.
  \end{enumerate}
  \begin{resposta}
  \end{resposta}

% ================================================================
%  BLOCO 3 — PROJETO E DECOMPOSIÇÃO
% ================================================================

  \vspace{0.4cm}

  \item \textbf{(Projeto -- CARRO)} Baseado no exemplo da aula, considere a relação:
  \texttt{CARRO(Placa, Modelo, Marca, Acessorio, Motorista)}.
  Sabe-se que:
  \begin{itemize}
    \item \texttt{Placa $\to$ \{Modelo, Marca\}}
    \item \texttt{Placa $\twoheadrightarrow$ Acessorio}
    \item \texttt{Placa $\twoheadrightarrow$ Motorista}
  \end{itemize}

  \begin{center}
  \begin{tikzpicture}[node distance=1cm and 2cm, font=\footnotesize\sffamily, thick]
      \node (p) [draw, fill=orange!20, circle] {Placa};
      \node (m) [draw, fill=gray!10, rectangle, above right=of p] {Modelo, Marca};
      \node (a) [draw, fill=blue!10, rectangle, right=of p] {Acessório};
      \node (mo) [draw, fill=green!10, rectangle, below right=of p] {Motorista};
      
      \draw[->] (p) -- (m) node[midway, left, sloped] {\tiny DF};
      \draw[->>] (p) -- (a) node[midway, above] {\tiny DMV};
      \draw[->>] (p) -- (mo) node[midway, left, sloped] {\tiny DMV};
  \end{tikzpicture}
  \end{center}

  Apresente o modelo relacional final na 4FN, identificando claramente as chaves primárias e estrangeiras. Explique por que manter todos em uma única tabela é ineficiente.
  \begin{resposta}
  \end{resposta}

  \vspace{0.4cm}

  \item \textbf{(Projeto -- TI)} Uma empresa de software registra seus projetos na tabela \texttt{PROJETOS\_TI(Cod\_Proj, Nome\_Proj, Tecnologia, Colaborador)}. 
  Regras: 
  \begin{itemize}
      \item Um projeto usa várias tecnologias (Python, SQL, etc).
      \item Um projeto tem vários colaboradores alocados.
      \item As tecnologias usadas no projeto não dependem de quem está alocado.
  \end{itemize}

  \begin{center}
  \begin{tikzpicture}[node distance=1.5cm, font=\footnotesize\sffamily, thick]
      \node (cp) [draw, fill=orange!10, rectangle] {Cod\_Proj};
      \node (np) [draw, fill=gray!5, rectangle, above=of cp] {Nome\_Proj};
      \node (t) [draw, fill=blue!5, rectangle, right=of cp] {Tecnologia};
      \node (co) [draw, fill=green!5, rectangle, left=of cp] {Colaborador};
      
      \draw[->] (cp) -- (np);
      \draw[->>] (cp) -- (t);
      \draw[->>] (cp) -- (co);
  \end{tikzpicture}
  \end{center}

  Monte uma instância (tabela com dados) que viole a 4FN com pelo menos 2 colaboradores e 2 tecnologias para o mesmo projeto, e mostre a solução normalizada.
  \begin{resposta}
  \end{resposta}

  \vspace{0.4cm}

  \item \textbf{(Projeto -- Clínica)} Normalize para a 4FN a relação:
  \texttt{PACIENTE(ID\_Pac, Nome, Doenca\_Cronica, Alergia)}. 
  Considere que um paciente pode ter várias doenças e várias alergias, e estas informações são fatos independentes sobre o paciente. Mostre como as anomalias de atualização são evitadas após a separação das tabelas.

  \begin{center}
  \begin{tikzpicture}[node distance=1.5cm, font=\footnotesize\sffamily, thick]
      \node (id) [draw, fill=orange!10, rectangle] {ID\_Pac};
      \node (n) [draw, fill=gray!5, rectangle, above=of id] {Nome};
      \node (d) [draw, fill=red!5, rectangle, right=of id] {Doença\_C};
      \node (a) [draw, fill=yellow!5, rectangle, left=of id] {Alergia};
      
      \draw[->] (id) -- (n);
      \draw[->>] (id) -- (d);
      \draw[->>] (id) -- (a);
  \end{tikzpicture}
  \end{center}
  \begin{resposta}
  \end{resposta}



  \item \textbf{Uma relação está na 4ª Forma Normal (4FN) se, e somente se:}
  \begin{enumerate}[label=\alph*)]
    \item Estiver na 3FN e não possuir dependências transitivas.
    \item Estiver na BCNF e toda dependência multivalorada não trivial tiver como determinante uma superchave.
    \item Estiver na 1FN e todos os atributos forem atômicos.
    \item Estiver na BCNF e não possuir dependências funcionais parciais.
    \item Não possuir chaves estrangeiras nulas.
  \end{enumerate}

  \item \textbf{A Dependência Multivalorada (DMV) $X \twoheadrightarrow Y$ é considerada não trivial quando:}
  \begin{enumerate}[label=\alph*)]
    \item $Y$ é um subconjunto de $X$.
    \item $X \cup Y$ é igual ao conjunto de todos os atributos da relação.
    \item $Y$ não é um subconjunto de $X$ e $X \cup Y$ não contém todos os atributos da relação.
    \item $X$ é a chave primária da relação.
    \item $Y$ depende funcionalmente de $X$.
  \end{enumerate}

  \item \textbf{O principal problema causado pela violação da 4FN é:}
  \begin{enumerate}[label=\alph*)]
    \item A perda de integridade referencial nas chaves estrangeiras.
    \item A redundância de dados gerada por um produto cartesiano entre atributos independentes.
    \item A impossibilidade de realizar junções (JOIN) entre tabelas.
    \item A criação de tabelas com excesso de chaves primárias.
    \item O aumento do desempenho das consultas de leitura.
  \end{enumerate}

  \item \textbf{Considere a relação $R(A, B, C)$ onde existem as DMVs $A \twoheadrightarrow B$ e $A \twoheadrightarrow C$. Se $B$ e $C$ são independentes, a decomposição correta para a 4FN gera as tabelas:}
  \begin{enumerate}[label=\alph*)]
    \item $R1(A, B)$ e $R2(B, C)$.
    \item $R1(A, B)$ e $R2(A, C)$.
    \item $R1(A, B, C)$ permanece inalterada.
    \item $R1(B, C)$ e $R2(A, B)$.
    \item $R1(A)$ e $R2(B, C)$.
  \end{enumerate}

  \item \textbf{Qual das alternativas abaixo melhor descreve a relação entre BCNF e 4FN?}
  \begin{enumerate}[label=\alph*)]
    \item BCNF é uma forma normal mais rigorosa que a 4FN.
    \item Uma relação na 4FN está obrigatoriamente na BCNF.
    \item Uma relação na BCNF está obrigatoriamente na 4FN.
    \item A 4FN trata de dependências funcionais, enquanto a BCNF trata de DMVs.
    \item Elas são conceitos independentes e não possuem relação hierárquica.
  \end{enumerate}

  \item \textbf{Em uma tabela que viola a 4FN, se um funcionário possui 3 telefones e 4 habilidades independentes, quantas linhas ele ocupará na tabela original?}
  \begin{enumerate}[label=\alph*)]
    \item 3 linhas.
    \item 4 linhas.
    \item 7 linhas.
    \item 12 linhas (produto cartesiano).
    \item 1 linha com campos multivalorados.
  \end{enumerate}



  \item \textbf{Uma DMV $X \twoheadrightarrow Y$ implica que, para um par de tuplas $(x, y1, z1)$ e $(x, y2, z2)$, também devem existir na relação as tuplas:}
  \begin{enumerate}[label=\alph*)]
    \item $(x, z1, y2)$ e $(x, z2, y1)$.
    \item $(x, y1, z2)$ e $(x, y2, z1)$.
    \item $(y1, x, z1)$ e $(y2, x, z2)$.
    \item $(x, x, x)$ e $(y, y, y)$.
    \item A DMV não exige a existência de outras tuplas.
  \end{enumerate}

  \item \textbf{Sobre as "Dependências Multivaloradas Independentes", assinale a correta:}
  \begin{enumerate}[label=\alph*)]
    \item Duas DMVs são independentes se o determinante de uma for o dependente da outra.
    \item Elas ocorrem quando dois atributos multivalorados não possuem ligação lógica entre si, mas dependem da mesma chave.
    \item Devem sempre ser mantidas na mesma tabela para facilitar o JOIN.
    \item São resolvidas através da 3FN.
    \item Não causam redundância se a tabela estiver na 1FN.
  \end{enumerate}




\end{enumerate}

\vspace{1cm}
\begin{center}
  \textit{``A 4FN remove redundâncias causadas por dependências multivaloradas, garantindo que fatos independentes não sejam forçados a coexistir em combinações artificiais.''}
\end{center}

\end{document}
